\documentclass{techbrief}

\usepackage{tikz, pgfplots}
\usetikzlibrary{positioning,calc,arrows,arrows.meta, shapes, decorations.markings}
\usetikzlibrary{decorations.pathreplacing}

\title{Physical vector spaces for quantum mechanics}
\author{Gabriele Carcassi}
\category{Quantum mechanics}
\tags{Hilbert spaces, Schwartz spaces, Functional analysis, Lie algebra}

\begin{document}
	
\maketitle
	
\begin{abstract}
	We develop a definition for a vector space that is more appropriate to respect physical requirements when representing quantum state spaces. Observables are continuous everywhere defined symmetric linear operators, which makes them better behaved under limits and change for reference frames. The Lie algebra becomes part of the definition of the space, which is what is already common practice in physics, to avoid mathematical equivalences that do not correspond to physical ones. The space is complete in the topology generated not just by the inner product, but also by the observables.
\end{abstract}

\section{Introduction}

As representation of quantum systems through Hilbert spaces fails to satisfy basic physical requirements, we look for an alternative that is both more physically meaningful and mathematically elegant. In the current form, we have developed a well-motivated and well-behaved structure for the topology, the observables and the inner product. The structure reproduces features of Schwartz space and recovers assumptions that most physicists expect but are not true in Hilbert spaces.

\section{Mathematical overview}

Instead of working with Hilbert spaces, we will give a characterization of a vector space that is better suited to capture physical requirements. It requires observables to be everywhere defined continuous self-adjoint operators, and ties the algebra of observable to the space definition. The space is also complete in the coarsest topology that makes both the observables and the inner product continuous.

\begin{defn}
	A \textbf{quantum vector space} is a complete second countable complex topological vector space $V$ that satisfies the following:
	\begin{itemize}
		\item $V$ is equipped with an inner product $\< \cdot , \cdot \> : V \times V \to \mathbb{C}$; an observable is a continuous everywhere defined self-adjoint linear operator $X : V \to V$
		\item $V$ is equipped with a set of defining observables $\{X_i\}_{i \in I}$ satisfying a given Lie algebra of commutation relationships $\frac{[X_i, X_j]}{\imath \hbar} = \sum\limits_k C_{ij}^k X_k$
		\item the topology is the coarsest vector space topology making the defining observables and the inner product continuous.
	\end{itemize}
\end{defn}

\begin{coro}
	The topology coincides with the one generated by the semi-norms of the form $\sqrt{\<Av, Av\>}$, where $A$ is an element of the unital algebra generated by the defining observables. Consequently, $V$ is a locally convex topological vector space.
\end{coro}

\begin{proof}
	Let $\mathcal{A}$ be the unital algebra generated by the defining observables. For each $A \in \mathcal{A}$, define the semi-norm $p_A(v) = \sqrt{\<Av, Av\>}$ with $v \in V$. These semi-norms generate a topology $\tau_{\mathcal{A}}$, which we want to show is exactly the topology of $V$.
	
	First of all, since $\tau_{\mathcal{A}}$ is generated by semi-norms, it is a locally convex vector space topology. Moreover, since $p_A(X_i v) = p_{A X_i}(v)$, and $A X_i \in \mathcal{A}$, $X_i$ is continuous. Lastly, since $p_I(v) = \Vert v \Vert$ is one of the semi-norms and $|\< v, w \>| \leq \Vert v \Vert \Vert w \Vert$ is dominated by semi-norms that generate the topology, the inner product is continuous. Therefore, $\tau_{\mathcal{A}}$ is a topology that makes the defining observables and the inner product continuous.
	
	For minimality, let $\tau$ be a vector space topology that makes the defining observables and the inner product continuous. Then $v \to \Vert v \Vert$ is $\tau$-continuous. Every element of $\mathcal{A}$ is $\tau$-continuous because the composition of continuous operator is a continuous operator. Therefore, $p_A(v) = \Vert A v \Vert$ is $\tau$-continuous for every $A \in \mathcal{A}$. So $\tau$ must contain $\tau_{\mathcal{A}}$. Therefore, $\tau_{\mathcal{A}}$ is the coarsest topology that makes the defining observables and the inner product continuous.
\end{proof}

\begin{example}[Schwartz]
	A Schwartz space $\mathcal{S}(\mathbb{R}^n)$ equipped with defining observables $\{X^i, P_i, I\}$ and algebra $\frac{[X^i, P_j]}{\imath \hbar} = \delta^i_j I$ is a quantum vector space. Moreover, two Schwartz spaces $\mathcal{S}(\mathbb{R}^n)$ and $\mathcal{S}(\mathbb{R}^m)$ are isomorphic as quantum vector spaces if and only if $n=m$.
\end{example}

\begin{remark}
	Note that the topology of the quantum vector space is exactly the topology of the Schwartz space as the $L^2$ semi-norms dominate the more standard sup norms and vice-versa. While Schwartz spaces of different dimensionality are isomorphic as topological vector spaces, they are not as quantum vector spaces because no linear transformation can change the number of defining observables.
\end{remark}

\begin{example}[Qubit]
	A two dimensional complex vector space equipped with defining observables $S^i = \frac{\hbar}{2}\sigma^i$, where $\sigma^i$ are the Pauli matrices, and algebra $\frac{[S^i, S^j]}{\imath \hbar} = \epsilon_{ijk} S^k$ is a quantum vector space.
\end{example}

\section{Conditions}

The following are conditions that we would expect any physical theory to obey, and therefore they should be considered part of a base theory for all of physics.

\begin{condition}[Continuity of measurable quantities]{[B]COBS}
	Small state changes lead to small changes of measurable quantities. That is, the expectation of a quantity is a continuous function of the state.
\end{condition}

\begin{condition}[Measurable quantities well-defined in all frames]{[B]WDQ}
	All states can equivalently be expressed in all frames. That is, if a quantity is well-defined in one frame, the corresponding quantity in another frame will be well-defined.
\end{condition}

\begin{condition}[Distinguishability of different systems]{[B]STDIST}
	Different physical systems have a different mathematical representation. That is, the mathematical representation of different physical systems must have enough structure to tell them apart.
\end{condition}

The following condition sums up the definition of quantum states on top of quantum vector spaces. We will test this against the above requirements.

\begin{condition}[Quantum states on a quantum vector space]{QSTQ}
	The state space of the system is the projective space of a quantum vector space $V$. That is, a state $\psi \in P(V)$ is represented by a ray $\psi = \{c |\psi\> \mid c \in \mathbb{C} \}$ for some $|\psi\> \in V$. The conditional probability $p : P(V) \times P(V) \to [0,1]$ is given by the Born rule. That is, $p(\psi|\phi) = \frac{\< \phi | \psi \>\< \psi | \phi \>}{\< \psi | \psi \>\< \phi | \phi \>}$.
\end{condition}

\section{Theorem}

Now we show that \ref{QSTQ} is consistent with each of the above base conditions.

\begin{thrm}
	Condition \ref{QSTQ} is consistent with \ref{[B]COBS}, \ref{[B]WDQ} and \ref{[B]STDIST}.
\end{thrm}

\begin{proof}
	Observables are defined to be continuous functions over $V$. Since the inner product is also continuous, the expectations of all observables is continuous. This means \ref{QSTQ} is consistent with \ref{[B]COBS}.
	
	Suppose we have two frames, $A$ and $B$, each with its own set of defining observables $X_i^A$ and $X_i^B$. Since the underlying system is the same, the defining observables of $A$ must be observables in $B$ and vice-versa, as each observer must be able to transform quantities to and from all frames. Since observables are everywhere-defined, there is no domain mismatch so the observables will be well-defined in all frames. This means \ref{QSTQ} is consistent with \ref{[B]WDQ}.
	
	Lastly, an isomorphism in the category of quantum vector spaces must preserve the Lie algebraic structure. The Lie algebraic structure exactly represents the observables and the symmetries that define the physics. Since two physically distinct systems must have physically distinct Lie algebras, two systems are physically equivalent if and only if their Lie algebras of observables are equivalent. This means \ref{QSTQ} is consistent with \ref{[B]STDIST}.
\end{proof}

\section{Additional remarks}

\textbf{Second countability vs separability.} It is standard to ask Hilbert spaces to be separable. For metrizable LCTVS, the two requirements are identical. For LCTVS in general, they are not. Second countability implies separability, but not vice-versa. From \href{https://assumptionsofphysics.org/autogen/AssumptionsOfPhysicsDraft.pdf#section.2.1.4}{Physical Mathematics}, we know that any set of physically definable objects is second countable in the natural topology. The most natural space to consider is the set of ensembles, which, if second countable, would force the set of pure states to be second countable, as they are a subspace of the first. If the pure states are second countable, we can always extend the topology to the vector space to preserve second countability.

\textbf{Number of defining observables.} For a field theory, we would expect the defining observable to be at least countable. Separability/second countability would allow us to find a countable subset of observables from a bigger set that generates the same topology.

\textbf{Completeness of the Lie algebra.} There is still a known problem with the current definition in that it allows spaces with the algebra that is, in some sense, not ``complete.'' For example, a Hilbert space with the defining observables of a qubit would be, with the current definition, a quantum vector space. A Schwartz space with only position in the defining observable would be, with the current definition, a quantum vector space. Both of these cases are pathological: the first has ``too many vectors,'' as it is the tensor product of the qubit and another space whose elements cannot be distinguished with the defining observables; the second has ``too few observables,'' as not all elements of the vector space on which the defining observables are defined can be distinguished with the defining observables. Addressing this gap is an open problem.

\textbf{Symmetric vs self-adjoint.} In the current structure, there is no difference between symmetric and self-adjoint, as all operators are defined on the whole space. In retrospect, this is the most natural solution from both the mathematical (i.e. in all other branches of math map are defined on the full space) and the physical (i.e. physicists seldom differentiate between the two properties) standpoint. What is lost is that not all self-adjoint operators generate a unitary transformation, but this is not as useful as it may seem. Symmetrized polynomials of position and momentum, in fact, are not always self-adjoint, and to understand which ones are the integrability of the system must be checked. This change brings quantum mechanics more in line with Hamiltonian mechanics, in which all functions are equivalent, though some do not generate a complete vector field.

\section{Conclusion}

The proposed quantum vector spaces capture basic physical requirements with minimal modification to what physicists are already used to. In fact, it better matches physical intuition. Once the open problem of the completeness of the Lie algebra is addressed, this note will be updated and we will have hopefully a final definition.

\end{document}

